\usepackage{amsmath}
\usepackage{amsfonts}
\usepackage{amssymb}
\usepackage{amsthm,mathabx}
\usepackage{graphicx}
\usepackage{geometry}
\usepackage{tikz}
\usepackage{tkz-tab} %graphiques et tableaux de variations
\usepackage[skins,breakable]{tcolorbox} % boîtes en environnements colorés
\usepackage{pgfplots}
\pgfplotsset{compat=1.18}

\usepackage{calc}
\usepackage{cancel}
\usepackage{multicol}
\usepackage[babel=true]{csquotes}
\usepackage{xspace}
\usepackage{setspace}
\usepackage{tabularx}
\usepackage{mdframed}
\usepackage{hyperref}
\usepackage{xcolor}
\usepackage{ifthen}
\usepackage{needspace}

\usepackage{qrcode}
\usepackage{mathrsfs}
%\usepackage{frcursive}

\usepackage{fancyvrb} %saveVerbatim
%\usepackage{verbatim} 

% Les environnements verbatim ne peuvent pas être lus dans l'argument de
% \contenu. \BVerbatimInput lit le fichier au moment où le contenu est affiché.
% Les extraits Python sont stockés dans code/python/ puis insérés par cette
% commande, qui reste utilisable à l'intérieur de \contenu{...}.
\newcommand{\pythoncode}[2][]{%
	\fbox{\BVerbatimInput[fontsize=\small,#1]{\exercisespath code/python/#2}}%
}

% Animation interactive d'OpenYourMath, \animation{nom}{légende} : le site ouvre
% la page animations/<nom>.html ; le PDF n'en garde que la légende, l'adresse et
% un QR code de l'adresse. La légende est du texte simple, sans formule.
% Même présentation que l'export LaTeX d'OpenYourMath (src/lib/latex/export.js).
\newcommand{\animation}[2]{%
	\par\medskip\noindent\fbox{\parbox{0.94\linewidth}{%
		\parbox[c]{\dimexpr\linewidth-2.4cm\relax}{\emph{Animation interactive : #2}\\
		\url{https://openyourmath.org/animation/#1}}\hfill
		\parbox[c]{2cm}{\qrcode[height=2cm]{https://openyourmath.org/animation/#1}}}}\par\medskip%
}

% Racine du dépôt, à définir avant le chargement des préambules si nécessaire.
% Ce nom évite la collision avec \path, commande TikZ.
\providecommand{\exercisespath}{}

%métadonnées :

\newcommand{\theme}[1]{}
\newcommand{\auteur}[1]{}
\newcommand{\organisation}[1]{}
\newcommand{\uuid}[1]{}
\newcommand{\datecreate}[1]{}
\newcommand{\similaire}[1]{}
\newcommand{\licence}[1]{}
\newcommand{\exoid}[1]{}
\newcommand{\video}[1]{}
\newcommand{\isIndication}[1]{}
\newcommand{\isCorrection}[1]{}
\newcommand{\chapitre}[1]{}
\newcommand{\sousChapitre}[1]{}
\newcommand{\niveau}[1]{}
\newcommand{\module}[1]{}
\newcommand{\difficulte}[1]{}
\newcommand{\cpge}[1]{} % rattachements aux programmes CPGE (référentiel d'exobase)
